\section{Literature context and source boundary}
\label{sec:literature-source}

\subsection{Classical ingredients and the paper-specific question}

The projective line over a ring can be defined through admissible or
unimodular pairs modulo unit scaling.  Its functorial geometry is developed
systematically by \citet{blunck2000projective}, while finite examples expose
how zero divisors change incidence and distantness
\citep{saniga2007classification}.  We use only the elementary commutative-ring
case $P^1(\mathbb Z/n\mathbb Z)$ and prove the required count directly.  The
state space itself is not a novelty claim.

The projective actions of order two and three belong to the modular-group and
Farey-graph setting.  Generalized Farey graphs encode modular actions on
finite quotients \citep{jones1991modular}, and symbolic codings for the
modular surface connect continued fractions, geodesic dynamics, and
cross-sections \citep{katok2006symbolic}.  Our use is deliberately
adversarial: universal modular relations are tested as a candidate
arithmetic recurrence mechanism, then shown to reproduce composite cycles.

Transfer operators for modular and Farey systems already produce
Fredholm-determinant descriptions of Selberg or Ruelle zeta functions.
\citet{mayer1991thermodynamic} gives the foundational modular-surface
construction; \citet{chang2001extension} treats general modular groups with
finite-dimensional representation data; and \citet{bonanno2014thermodynamic}
develops a two-variable Farey-map formalism.  These are the closest analytic
ancestors.  We do not identify our determinant with a Selberg, Ruelle, or
Riemann zeta function.  Instead, we ask whether the all-modulus graph-step
operator itself belongs to the trace class and what arithmetic classes its
owned traces contain.

\begin{table}[t]
\centering
\caption{Classical ingredients versus the paper-specific use and claim
boundary.}
\label{tab:literature-boundary}
\small
\begin{tabularx}{\textwidth}{L{0.23\textwidth}YY}
\toprule
Ingredient & Paper-specific use & Claim boundary \\
\midrule
Projective lines over rings & Source-derived state spaces $X_n$ and static
field defect & State space and count are classical \\
Modular/Farey actions & Nonterminal shared-state recurrence and generic
presentation control & Relations $S^2=R^3=1$ are classical \\
Transfer determinants & Same-object trace-class test for the frozen
all-modulus graph & No Selberg/Ruelle/Riemann identity is claimed \\
Trace-class determinants & Ordinary meaning of $\det(I-zB_s)$ & No
regularized or formal product is substituted \\
\bottomrule
\end{tabularx}
\end{table}

Within the documented search through August 15, 2026, we found no directly
comparable paper combining the frozen semiring-source grammar, a complete
preweight composite-cycle audit, and ordinary Fredholm ownership of the same
uninduced object.  This is a search-bounded positioning statement, not an
absolute priority claim.  Chang--Mayer is the nearest analytic collision;
the present contribution is the controlled obstruction, not the classical
modular machinery.

\subsection{Finite-full-shift source}

Let
\[
 F_n=(A_n^{\mathbb Z},\sigma_n),\qquad n\in\Nzero,
\]
denote the conjugacy class of the full shift on an $n$-symbol alphabet.  The
source retains the operations
\[
 F_m\boxplusop F_n\cong F_{m+n},\qquad
 F_m\boxtimesop F_n\cong F_{mn},
\]
the zero $F_0$, unit $F_1$, successor, source equality,
quotient/remainder, congruence, and entropy $h(F_n)=\log n$ for $n\ge1$.
Alphabet sum is an operation on full-shift classes; no categorical coproduct
claim is required.

For each $n\ge2$, source congruence reconstructs
$\Rring_n=\Z/n\Z$.  Units, additive inverse, unimodular pairs, and projective
classes are defined by source equations.  Constants 2 and 3 are generated
from the unit, and $[1:0]$ is a source-defined projective point.  Consequently
the state and edge formulas below are invariant under any isomorphism that
transports the complete semiring/congruence presentation.

\begin{convention}[Source-natural construction]
\label{conv:source-natural}
A candidate is source-natural only if an isomorphic relabeling transports
every operation, relation, state, edge, roof, and marker.  Exact matched-clone
agreement is required and is not counted as evidence of arithmetic
selectivity.
\end{convention}

\subsection{Information boundary and evaluation order}

The candidate may use the source operations, exact finite-ring arithmetic
that implements them, the two fixed matrices, cusp, frozen roofs, deterministic
cutoffs and seeds, and evaluator-only labels applied after enumeration.  It
may not use a supplied prime, factor, prime-power, accepted-support,
orbit-projector, or target-zero table.  Candidate-side primality or
factorization calls, fitted coefficients, Riemann-zero ordinates, root
matching, and Route B are excluded.

The evaluation order is load-bearing.  We first freeze the transition graph,
then prove its complete primitive support, then run clone and random-action
controls, and only afterward study operator class and determinants.  The
static equality $|X_n|=n+1$ may be reported as a source fact, but it may not
decide which blocks or edges exist.  This order prevents an exact primality
test from being renamed as emergent periodic structure.
